\documentclass[10pt,journal,compsoc]{IEEEtran}
\usepackage{cite}
\usepackage{amsmath,amssymb,amsfonts,amsthm,mathtools}
\usepackage{booktabs}
\usepackage{enumitem}
\usepackage{url}
\usepackage{xr}
\externaldocument[main-]{build/paper1}

\newtheorem{theorem}{Theorem}
\renewcommand{\thetheorem}{S.\arabic{theorem}}
\newtheorem{lemma}[theorem]{Lemma}
\newtheorem{corollary}[theorem]{Corollary}
\newtheorem{proposition}[theorem]{Proposition}
\newtheorem{definition}[theorem]{Definition}
\renewcommand{\theequation}{S.\arabic{equation}}

\newcommand{\Nat}{\mathbb{N}}
\newcommand{\Real}{\mathbb{R}}
\newcommand{\satM}{\models_{\mathrm{MTL}}}
\newcommand{\satL}{\models_{\mathrm{LTL}}}
\newcommand{\Until}{\mathbin{\mathsf{U}}}
\newcommand{\Release}{\mathbin{\mathsf{R}}}
\newcommand{\hUntil}{\mathbin{\widehat{\mathsf{U}}}}
\newcommand{\hRelease}{\mathbin{\widehat{\mathsf{R}}}}
\newcommand{\Next}{\mathsf{X}}
\newcommand{\Always}{\mathsf{G}}
\newcommand{\Eventually}{\mathsf{F}}
\newcommand{\WeakUntil}{\mathbin{\mathsf{W}}}
\newcommand{\Reset}{\mathsf{rst}}
\newcommand{\Unch}{\mathsf{unch}}
\newcommand{\T}{\mathcal{T}}
\newcommand{\CVal}{\mathsf{val}}
\newcommand{\CReset}{\mathsf{reset}}
\newcommand{\elapsed}{\mathsf{elap}}
\newcommand{\ewbase}{\mathsf{ew\_base}}
\newcommand{\Node}{\mathsf{node}}
\newcommand{\runinhead}[1]{\par\smallskip\noindent\textbf{#1}\enspace}

\title{Supplementary Material for ``A Mechanically Verified Translation of MTL$_{0,\infty}$ into Timed B\"uchi Automata''}
\author{Elie Fares, Jean-Paul Bodeveix, Hussain Al-Aqrabi, and Azar Salami}
\begin{document}
\maketitle

This supplement contains the detailed reset-policy and residual-obligation
arguments, reset-completion proofs, and auxiliary clock arithmetic. The
Rocq source files used by the development are included in the accompanying
\texttt{proof/} directory. The complete project and build instructions are at
\url{https://github.com/efares1/coq-verified-mtl-to-tba}. Cross-references
point to the main paper.

\noindent\textbf{Proof map.} The arguments below support the four challenges
in main-paper Section~\ref{main-subsec:mechanization-challenges}: occurrence-
indexed induction in Supplement Section~\ref{subsec:structural-equivalence};
repeated-activation completeness in Section~\ref{subsec:timed-completeness};
relaxation and reset completion in Sections~\ref{subsec:reset-completion-proof}
and~\ref{subsec:relaxation-exactness}; and strict-bound reasoning in the
primitive arguments of Sections~\ref{subsec:timed-soundness} and
\ref{subsec:timed-completeness}. The Rocq development contains the complete
machine-checked proofs; this supplement explains their main invariants.

\section{Omitted Proofs for the TBA Construction}
\label{sec:supp-proofs}

\subsection{Reset Completion}
\label{subsec:reset-completion-proof}
Proof of Lemma~\ref{main-lem:reset-completion}. The extension $\rho'$ is
clock-consistent by construction. Consider a preservation clock $x$. If the
cube taken at position $i$ contains $\Unch(x)$, then neither $\rho$ nor
$\rho'$ resets $x$ at $i$; otherwise $\rho'$ resets it. Hence $\rho'$ resets
$x$ wherever $\rho$ does, and the clock recurrence
\eqref{main-eq:clock-update} gives, by induction on positions,
$\CVal_{\rho'}(i,x)\le\CVal_\rho(i,x)$. The literals over $x$ are upper
bounds, so remain true. Dually, $\rho'$ resets a restart clock only where
the cube contains $\Reset(x)$, hence only where $\rho$ resets it. Therefore
$\CVal_{\rho'}(i,x)\ge\CVal_\rho(i,x)$, and the lower-bound literals
remain true. Action literals and the B\"uchi condition are unchanged, so the
same run is accepting after completion (Rocq lemma
\texttt{complete\_complete}).

\subsection{Exactness of Relaxation and Reset Completion}
\label{subsec:relaxation-exactness}
Proof of Proposition~\ref{main-prop:reset-completion}. For soundness, a run of the
completed TBA gives a run of the relaxed automaton on the propositional word
of its clock-consistent extension. At each position, the relaxed transition
comes from an original backend cube. Reconstruct a valuation satisfying that
cube by setting each relevant extended atom false when its negative literal
was deleted, retaining values required by surviving literals, and choosing
values for deleted irrelevant atoms as needed. The resulting word follows the
same state sequence in the original backend automaton, so the backend contract
makes it a model of $\T(f)$. The reconstructed and relaxed words agree on base
atoms; on relevant extended atoms, relaxation can only change false to true.
Since $\T(f)$ contains extended atoms only positively, it is upward closed in
those atoms, so the relaxed word also satisfies $\T(f)$. The soundness
direction of Theorem~\ref{main-thm:encoding-correctness} gives
$w,0\satM f$.

For completeness, Theorem~\ref{main-thm:encoding-correctness} supplies a
clock-consistent extension of $w$ satisfying $\T(f)$. The backend contract
therefore gives an accepting run of the original backend automaton.
Relaxation preserves this run, and the reset-completion lemma supplies an
extension accepted by the completed TBA. The mechanized theorem is
\texttt{MTL\_to\_TBA\_correct\_with}.

\section{Detailed Proofs of Soundness and Completeness}
\label{sec:supp-proof}
The main paper states the semantic model, translation clauses, residual predicates, canonical reset policy, and proof roadmap. This supplement gives the omitted primitive soundness and completeness arguments, the structural induction, and the detailed theorem composition. It uses the main-paper definitions rather than restating them.

\subsection{Soundness of the Primitive Timed Translations}
\label{subsec:timed-soundness}

For each primitive hatted timed constructor $F$, the local soundness lemma has
the form
\begin{equation}
  \rho,i\satL \T_\pi(F)
  \quad\Longrightarrow\quad
  \ewbase(\rho),i\satM F,
\label{eq:local-soundness}
\end{equation}
assuming clock consistency and the induction hypotheses for the immediate
subformulas.  No canonical reset-policy assumption appears in
\eqref{eq:local-soundness}.

\runinhead{Upper-bounded hatted Until.}
The translated LTL Until starts at $i+1$ and provides an endpoint $j>i$.
At position $i$ itself the translation does not constrain the clock, but the
recurrence~\eqref{main-eq:clock-update} together with nonnegativity gives
$\CVal_\rho(i+1,x)\ge\Delta_i$ whether or not $x$ is reset at $i$.  Along the
intermediate positions $(i,j)$ the shared clock is unchanged, so
Lemma~\ref{main-lem:clock-accumulation} applied on $[i+1,j)$ yields
$\CVal_\rho(j,x)\ge\Delta_i+\elapsed_w(i+1,j)=\elapsed_w(i,j)$, and the
endpoint guard $x\le d$ gives
\[
  \elapsed_w(i,j)\le d.
\]
The induction hypotheses turn the translated prefix $\overline{P}$ into $p$ and the
endpoint $\overline{Q}$ into $q$, matching the semantics of
$p\hUntil_{\le d}q$.  For $\hUntil_{<d}$, the strict guard $x<d$ gives
$\elapsed_w(i,j)<d$ in the same way (\texttt{sound\_\allowbreak{}UhatLt}).

\runinhead{Lower-bounded hatted Until.}
The translation resets $x$ at $i$ and follows $\Gamma$.  In the finite branch,
a future state satisfies $x\ge d$ together with $\overline{P}\Until \overline{Q}$.  By
\eqref{main-eq:lower-clock}, the endpoint is at least $d$ time units after $i$, and
the nested Until supplies the required $p$-prefix and $q$-endpoint.  In the
second branch, $\Always \overline{P}\land\Always\Eventually \overline{Q}$ holds.  Then $p$ holds
forever and $q$ recurs infinitely often; time divergence guarantees a $q$
occurrence at physical distance at least $d$.  For $\hUntil_{>d}$, the strict
threshold $x>d$ gives a strict lower bound in the finite branch, and the
fairness branch uses the strict argument of Section~\ref{main-subsec:main-strict-bounds}
(\texttt{sound\_\allowbreak{}UhatGt}).

\runinhead{Upper-bounded hatted Release.}
The translation resets the shared clock and requires
\[
  \overline{Q}\WeakUntil\bigl((x>d)\lor(\overline{P}\land \overline{Q})\bigr).
\]
Before the weak-until endpoint, $q$ holds at every strictly future position.
If the endpoint is $\overline{P}\land \overline{Q}$, $p$ discharges the Release obligation; if it is
$x>d$, the upper-bounded window has ended.  For $\hRelease_{<d}$, the window
ends at $x\ge d$ (\texttt{sound\_\allowbreak{}RhatLt}).

\runinhead{Lower-bounded hatted Release.}
The translation requires
\[
  ((x<d)\land\Unch(x))
  \WeakUntil
  \bigl((\overline{P}\Release \overline{Q})\lor((x<d)\land \overline{P})\bigr)
\]
from $i+1$.  While the left side holds, the clock is preserved and remains
below $d$, so those positions are still before the lower threshold.  At the
endpoint either an ordinary untimed Release takes over or $p$ discharges the
obligation before the threshold is reached.  For $\hRelease_{>d}$, the
pre-threshold test is $x\le d$ (\texttt{sound\_\allowbreak{}RhatGt}).

There are no separate ordinary timed soundness lemmas.  After unfolding
Eqs.~\eqref{main-eq:derive-ule}--\eqref{main-eq:derive-rge}, their
correctness follows from the Boolean cases of the structural induction and the
eight hatted lemmas above.

\subsection{Completeness under the Canonical Reset Policy}
\label{subsec:timed-completeness}

The following case arguments prove the converse on the canonical extension
defined in the main paper:
\[
  w,i\satM F
  \quad\Longrightarrow\quad
  \rho^\star,i\satL\T_\pi(F).
\]
They use the reset decisions, clock update, and residual predicates already
given in Table~\ref{main-tab:canonical-policy} and
Eq.~\ref{main-eq:clock-update}; the details below establish how each primitive
case satisfies its translated obligation.

\runinhead{Upper-bounded hatted Until.}
Choose the first $q$-position among the MTL witnesses.  The proof maintains,
for every position $n$ with $i<n\le j$ up to that endpoint $j$, the invariant
\[
  \CVal_{\rho^\star}(n,x)+\elapsed_w(n,j)\le d.
\]
It need not hold at the activation position $i$ itself (the policy may reset
there because the old value is useless); it is established at $i+1$ in both
the reset and the preservation case.
If the policy preserves the clock, Table~\ref{main-tab:canonical-policy}
guarantees an older reference whose residual $\mathsf{URes}$ is still
satisfiable strictly after the current position; if it resets, the new value is
consistent with the current activation.  Hence every intermediate position satisfies the clock
guard required by the translated Until and the first $q$-position supplies its
endpoint.  This is the content of the mechanized lemma
\texttt{complete\_\allowbreak{}UhatLe}.  The strict case maintains the strict
invariant $\CVal_{\rho^\star}(n,x)+\elapsed_w(n,j)<d$ with the residual
$\mathsf{URes}^{<}$ (\texttt{complete\_\allowbreak{}UhatLt}).

\runinhead{Lower-bounded hatted Until.}
At every position where $F$ holds, the canonical policy resets the shared
clock (Table~\ref{main-tab:canonical-policy}).  The proof propagates an active obligation until a position satisfies
\[
  x\ge d
  \quad\text{and}\quad
  \overline{P}\Until \overline{Q},
\]
which yields the finite branch of $\Gamma$.  If no such position is reached,
the active-obligation invariant implies that $p$ remains continuously true,
and since every later position again carries a pending obligation that
requires a later $q$, the proposition $q$ recurs infinitely often.  This gives
\[
  \Always \overline{P}\land\Always\Eventually \overline{Q},
\]
the second branch of $\Gamma$.  This is formalized by
\texttt{complete\_\allowbreak{}UhatGe}, and by
\texttt{complete\_\allowbreak{}UhatGt} for the strict threshold $x>d$.

\runinhead{Upper-bounded hatted Release.}
The canonical policy resets the clock at every position where $F$ holds.  The
proof follows the currently active nontrivial Release obligation.  Until it
is discharged, the clock measures elapsed time from the relevant start and
$q$ continues to hold.  The weak until terminates either because $p\land q$
discharges the Release or because $x>d$, meaning that the upper-bounded window
has ended.  This is the role of \texttt{complete\_\allowbreak{}RhatLe}, and of
\texttt{complete\_\allowbreak{}RhatLt} for the strict window end $x\ge d$.

\runinhead{Lower-bounded hatted Release.}
Here the invariant is
\[
  \mathsf{RRes}
  \bigl(n,\CVal_{\rho^\star}(n,x),d,p,q\bigr).
\]
By Table~\ref{main-tab:canonical-policy}, the shared clock is preserved exactly
while $x<d$, the residual holds, and $p$ does not occur; in that case the
residual shifts
from $(i,v)$ to
\[
  (i+1,v+\Delta_i).
\]
The Rocq lemma \texttt{rge\_\allowbreak{}residual\_\allowbreak{}shift} establishes this one-step update.
A converse lemma on the preservation policy recovers the residual information needed
at the next position, so the left side of the weak until remains valid until
either an ordinary untimed Release is established or $p$ occurs while the
clock is still below $d$.  This yields \texttt{complete\_\allowbreak{}RhatGe}.
The strict case uses $\mathsf{RRes}^{>}$, the pre-threshold test $x\le d$, and
the shift lemma \texttt{rgt\_\allowbreak{}residual\_\allowbreak{}shift}
(\texttt{complete\_\allowbreak{}RhatGt}).

As with soundness, ordinary timed operators need no additional completeness
lemmas: unfolding their definitions reduces them to Boolean structure and one
of the eight hatted primitive cases above.

\subsection{Structural Equivalence at Every Syntactic Occurrence}
\label{subsec:structural-equivalence}

The local timed lemmas are assembled with the homomorphic Boolean and untimed temporal
cases into the following stronger induction theorem.

\begin{theorem}[Canonical correctness at every occurrence]
\label{thm:canonical-all}
Let $\mathit{root}$ be well formed and suppose
\[
  \Node(\mathit{root},\pi)=f.
\]
Then for every position $i$,
\[
  w,i\satM f
  \quad\Longleftrightarrow\quad
  \rho^\star,i\satL\T_\pi(f),
\]
where $\rho^\star=\mathsf{canonical\_ext}(\mathit{root},w)$.
\end{theorem}

\begin{proof}[Proof idea]
The proof is by structural induction on the formula $f$.  Atomic propositions,
Boolean connectives, Next, untimed Until, and untimed Release follow directly
from their semantics and the induction hypotheses.  Since the ordinary timed
operators are Rocq definitions, they have already unfolded into these structural
cases.  The only primitive timed branches are therefore the eight hatted
constructors.  In each such branch, the forward direction invokes the
corresponding completeness lemma from
Section~\ref{subsec:timed-completeness}, while the reverse direction invokes
the corresponding soundness lemma from
Section~\ref{subsec:timed-soundness}.  The $\Node$ lemmas locate the two
children, and well-formedness of the located formula provides the bound and
recursive hypotheses.
\end{proof}

This is stronger than a root-only statement and is important in the mechanized
proof.  The induction retains the path in order to locate the current subformula
and its children, but the path is not the clock key.  If two paths locate the
same primitive timed formula, both branches refer to the same formula-keyed
clock.

\subsection{Proof of the Correctness Theorem}
\label{subsec:proof-main-theorem}
We now prove Theorem~\ref{main-thm:encoding-correctness}.

\begin{proof}
We prove the two implications separately.

\smallskip
\noindent\emph{($\Rightarrow$) MTL to clocked LTL.}
Assume $w,0\satM f$.  Choose the canonical extension
\[
  \rho=\mathsf{canonical\_ext}(f,w).
\]
By construction, $\rho$ has the same base timed word as $w$ and is
clock-consistent. Applying Theorem~\ref{thm:canonical-all} at the root path $[]$ and
position $0$ gives
\[
  \rho,0\satL\T(f).
\]
Therefore the required extension exists.

\smallskip
\noindent\emph{($\Leftarrow$) Clocked LTL to MTL.}
Assume there exists an arbitrary $\rho$ such that
\[
\begin{gathered}
  \mathsf{same\_base}(\rho,w),\qquad
  \mathsf{clock\_consistent}(\rho),\\
  \rho,0\satL\T(f).
\end{gathered}
\]
This extension need not be canonical, so completeness cannot be used.  Instead, a
structural induction on $f$ uses the policy-independent soundness lemmas of
Section~\ref{subsec:timed-soundness} and obtains
\[
  \ewbase(\rho),0\satM f.
\]
Finally, $\mathsf{same\_base}(\rho,w)$ rewrites the underlying timed word
$\ewbase(\rho)$ to $w$, yielding $w,0\satM f$.
\end{proof}


\section{Rocq Result Map}
\label{app:rocq}
\begingroup
\raggedright
The paper's principal statements correspond to these definitions and
theorems in the included Rocq development.\par\smallskip
\noindent\textbf{Proof source.} The Rocq definitions of
\texttt{EncodingCorrect\_proved} and
\texttt{MTL\_to\_TBA\_correct\_with} are in
\path{proof/EncodingCorrect_Shared_Clock_Derived_Strict_Direct_Proof.v}.\par\smallskip
\noindent\textbf{Core source.} The definitions of
\texttt{T\_uses\_every\_clock} and \texttt{complete\_complete} are in
\path{proof/MTL_to_TBA_Shared_Clock_Derived_Strict_Direct_Core.v}.\par\smallskip
The core file also contains the four ordinary Until and Release compatibility
lemmas and their strict variants.\par\smallskip
The path-indexed theorem \texttt{canonical\_encoding\_all} and the occurrence
lemmas \texttt{identical\_nodes\_share\_clock} and
\texttt{identical\_nodes\_share\_canonical\_trace} are in the correctness
module. The same module defines \texttt{ule\_residual},
\texttt{ult\_residual}, \texttt{rge\_residual}, and
\texttt{rgt\_residual}, together with the eight primitive soundness lemmas
and eight canonical-completeness lemmas. In the core module,
\texttt{psat\_mono}, \texttt{relax\_sound}, \texttt{relax\_complete}, and
the reset-completion lemmas discharge the backend-interface obligations.\par\smallskip
The theorem with suffix \texttt{\_with} assumes the backend correctness
contract as an argument and has no project-specific backend axiom. The standard
library assumptions are preserved in \path{proof/ASSUMPTIONS.txt}; reproduction
instructions are in \path{proof/README_proof.md}. The convenience theorem
\texttt{MTL\_to\_TBA\_correct} uses the core axiom
\texttt{LTL\_TO\_BUCHI\_CORRECT}; the main paper states the contract-
parameterized theorem.
\par
\endgroup

\section{Concrete Alarm TBA Listing}
\label{app:alarm-tba}
For the alarm instance in Section~\ref{main-sec:alarm-example} of the main paper,
the compiled TBA \texttt{alarm\_response\_spot\_tba} has locations
$\{0,1,2\}$, initial location $0$, accepting locations $\{0,1\}$, and clock
$x$. The following transitions result from applying relaxation and reset
completion to the manual Rocq transcription of the Spot graph. Here $A$ and
$H$ denote the alarm and handled events; $\top$ means no clock guard, and
$\varnothing$ means no reset. Guards are checked before resets.
\[
\begin{array}{c|c|c|c|c}
\text{source}&\text{target}&\text{event}&\text{guard}&\text{reset}\\ \hline
0&0&H&\top&\{x\}\\
0&0&\neg A&\top&\{x\}\\
0&1&A\land\neg H&\top&\{x\}\\
1&0&H&x\le 3&\{x\}\\
1&2&\neg H&x\le 3&\varnothing\\
2&0&H&x\le 3&\{x\}\\
2&2&\neg H&x\le 3&\varnothing
\end{array}
\]
The manual DOT-to-record correspondence is not verified; this table lists the
TBA compiled from the Rocq record whose propositional language contract is
proved in the main paper.

\end{document}
